\documentclass{article}
\usepackage[utf8]{inputenc}
\usepackage{amsmath} 
\usepackage{geometry}
\usepackage{listings}
\usepackage{xcolor}
\usepackage{array}

\geometry{a4paper, margin=1in}

\title{Labwork 1: Gradient Descent Analysis}
\author{Nguyễn Duy Anh Tùng - 2540107}
\date{\May 6, 2026}

\begin{document}

\maketitle

\section{Introduction}
This report describes the implementation of the Gradient Descent algorithm from scratch to find the minimum value of the function $f(x) = x^2$ with its derivative $f'(x) = 2x$.

\section{Implementation}
The algorithm follows these iterative steps:
\begin{enumerate}
    \item Initialize $x$ with a random value $x_0 = 10$.
    \item Update $x$ using the formula: $x = x - r \cdot f'(x)$, where $r$ is the learning rate.
    \item Repeat for a fixed number of epochs.
\end{enumerate}

\section{Experimental Results}
Below are the iteration steps for $r = 0.1$:

\begin{center}
    \begin{tabular}{|l|} 
        \hline 
        \begin{minipage}{0.5\textwidth} 
            \vspace{2mm}
            \lstinputlisting{result.txt}
            \vspace{1mm}
        \end{minipage} \\
        \hline 
    \end{tabular}
\end{center}

\section{Learning Rate Analysis}
Based on my experiments, the choice of learning rate $r$ significantly affects convergence:

\begin{itemize}
    \item \textbf{Small $r$ (e.g., 0.01):} The model converges very slowly. After 10 steps, $x$ is still far from 0.
    \item \textbf{Optimal $r$ (e.g., 0.1):} The model reaches the minimum point $x \approx 0$ quickly and smoothly.
    \item \textbf{Large $r$ (e.g., 1.1):} The updates are too drastic, causing $x$ to overshoot and diverge (values of $f(x)$ increase instead of decreasing).
\end{itemize}

\section{Conclusion}
Gradient Descent is highly sensitive to the learning rate. A "just right" learning rate is crucial for efficient training.

\end{document}
